
\documentclass{article}
\usepackage{colortbl}
\usepackage{makecell}
\usepackage{multirow}
\usepackage{supertabular}

\begin{document}

\newcounter{utterance}

\centering \large Interaction Transcript for game `codenames', experiment `risk\_high', episode 0 with mistral{-}large{-}3{-}azure.
\vspace{24pt}

{ \footnotesize  \setcounter{utterance}{1}
\setlength{\tabcolsep}{0pt}
\begin{supertabular}{c@{$\;$}|p{.15\linewidth}@{}p{.15\linewidth}p{.15\linewidth}p{.15\linewidth}p{.15\linewidth}p{.15\linewidth}}
    \# & $\;$A & \multicolumn{4}{c}{Game Master} & $\;\:$B\\
    \hline

    \theutterance \stepcounter{utterance}  
    & & & \multicolumn{2}{p{0.3\linewidth}}{
        \cellcolor[rgb]{0.9,0.9,0.9}{
            \makecell[{{p{\linewidth}}}]{
                \texttt{\tiny{[GM$|$GM]}}
                {'hidden': {'team': ['vinis', 'švinas', 'kranas', 'vašingtonas', 'kivis', 'kryžius', 'raketė', 'lenta', 'dramblio kaulas'], 'innocent': ['kirminas', 'masė', 'vaikas'], 'opponent': ['skylė', 'diena', 'bermudai', 'skristi', 'feniksas', 'suknelė', 'kentauras', 'žolė'], 'assassin': ['laboratorija', 'meksika', 'pastaba', 'atsargos', 'gyvenimas']}, 'revealed': {'team': {'team': [], 'innocent': [], 'opponent': [], 'assassin': []}, 'opponent': {'team': [], 'innocent': [], 'opponent': [], 'assassin': []}}}
            }
        }
    }
    & & \\ \\

    \theutterance \stepcounter{utterance}  
    & & \multicolumn{4}{p{0.6\linewidth}}{
        \cellcolor[rgb]{0.9,0.9,0.9}{
            \makecell[{{p{\linewidth}}}]{
                \texttt{\tiny{[P1$\langle$GM]}}
                \texttt{Pateikite vieną vieno žodžio užuominą, susijusią su kai kuriais jūsų komandos žodžiais iš toliau pateikto sąrašo. Galite pasirinkti, su kiek ir kokiais žodžiais bus susijusi jūsų užuomina.} \\
\texttt{Tikslas {-} kad jūsų partneris, remdamasis jūsų pateiktomis užuominomis, atspėtų visus jūsų komandos žodžius, kol kita komanda neatskleidė visų savo žodžių. Jei jūsų partneris atspėja "priešininko" arba "atitraukiančio dėmesį" žodį, jo spėjimo eilė baigiasi ir priešininkų komanda gali atlikti savo ėjimą. Jei jūsų partneris atspėja "žudantį" žodį, žaidimas baigiasi ir jūsų komanda iš karto pralaimi.} \\
\texttt{Kuo daugiau žodžių, su kuriais susijusi jūsų užuomina, tuo greičiau galite laimėti žaidimą, tačiau partneriui gali būti sunkiau atspėti asociacijas, todėl išmintingai pasirinkite užuominą. Užuominos žodis turi būti semantiškai susijęs su tiksliniais žodžiais, jis negali būti vienas iš sąrašuose esančių žodžių arba sudaryti jų dalių.} \\
\\ 
\texttt{Visada pateikite vieno žodžio užuominą ir kableliais atskirtą susijusių tikslinių žodžių sąrašą toliau nurodyta forma ir atsakymus pateikite kuo trumpesnius, niekada nepateikite jokio kito teksto, nei reikalaujama šioje formoje:} \\
\\ 
\texttt{UŽUOMINA: <ŽODIS>} \\
\texttt{TIKSLAI: <TIKSLAI>} \\
\\ 
\texttt{Jūsų komandos žodžiai yra: vinis, švinas, kranas, vašingtonas, kivis, kryžius, raketė, lenta, dramblio kaulas.} \\
\texttt{Jūsų priešininko komandos žodžiai yra: skylė, diena, bermudai, skristi, feniksas, suknelė, kentauras, žolė.} \\
\texttt{Atitraukimo žodžiai yra: kirminas, masė, vaikas.} \\
\texttt{Žudymo žodžiai yra: laboratorija, meksika, pastaba, atsargos, gyvenimas.} \\
            }
        }
    }
    & \\ \\

    \theutterance \stepcounter{utterance}  
    & \multicolumn{4}{p{0.6\linewidth}}{
        \cellcolor[rgb]{0.9,0.9,0.9}{
            \makecell[{{p{\linewidth}}}]{
                \texttt{\tiny{[P1$\rangle$GM]}}
                \texttt{UŽUOMINA: statyba} \\
\texttt{TIKSLAI: vinis, kranas, lenta} \\
\\ 
\texttt{{-}{-}{-}} \\
\texttt{UŽUOMINA: sunkus} \\
\texttt{TIKSLAI: švinas, dramblio kaulas} \\
\\ 
\texttt{{-}{-}{-}} \\
\texttt{UŽUOMINA: simbolis} \\
\texttt{TIKSLAI: kryžius, vašingtonas} \\
\\ 
\texttt{{-}{-}{-}} \\
\texttt{UŽUOMINA: kosmosas} \\
\texttt{TIKSLAI: raketė} \\
\\ 
\texttt{{-}{-}{-}} \\
\texttt{UŽUOMINA: vaisius} \\
\texttt{TIKSLAI: kivis} \\
            }
        }
    }
    & & \\ \\

    \theutterance \stepcounter{utterance}  
    & & & \multicolumn{2}{p{0.3\linewidth}}{
        \cellcolor[rgb]{0.9,0.9,0.9}{
            \makecell[{{p{\linewidth}}}]{
                \texttt{\tiny{[GM$|$GM]}}
                {'player': 'Cluegiver', 'type': 'rambling error', 'utterance': 'UŽUOMINA: statyba\nTIKSLAI: vinis, kranas, lenta\n\n---\nUŽUOMINA: sunkus\nTIKSLAI: švinas, dramblio kaulas\n\n---\nUŽUOMINA: simbolis\nTIKSLAI: kryžius, vašingtonas\n\n---\nUŽUOMINA: kosmosas\nTIKSLAI: raketė\n\n---\nUŽUOMINA: vaisius\nTIKSLAI: kivis', 'message': 'Your answer contained more than two lines, please only give one clue and your targets on two separate lines.'}
            }
        }
    }
    & & \\ \\

\end{supertabular}
}

\end{document}
